% CLAUDE REPORT::WCDERIV-0001::2026-11-12T01:51Z
% !TEX program = xelatex
% Compiles with XeLaTeX or pdfLaTeX (ASCII source).
\documentclass[11pt]{article}
\usepackage[margin=1in]{geometry}
\usepackage{amsmath,amssymb,amsthm}
\usepackage{booktabs,array,longtable}
\usepackage[hidelinks]{hyperref}
\usepackage{xurl}
\setlength{\emergencystretch}{3em}
\newtheorem{theorem}{Theorem}
\newtheorem{lemma}[theorem]{Lemma}
\newtheorem{corollary}[theorem]{Corollary}
\theoremstyle{definition}
\newtheorem{definition}[theorem]{Definition}
\newcommand{\Wb}{\overline W}
\newcommand{\Rm}{R_{\mathrm{mix341}}}
\newcommand{\poly}{\operatorname{poly}}
\newcommand{\cost}{\operatorname{cost}}
\title{Worst-Case Minimum Derivation Cost: Strength Calibration of the\\
Fixed-Calculus Lower Bound and the B3 Bridge (WCDERIV-0001)}
\author{Claude research agent (scheduled run, shared worst-case directive)}
\date{11 November 2026, 17:51 PST}
\begin{document}
\maketitle

\section*{Status}
\textbf{P versus NP remains OPEN.} Every load-bearing arrow of the
worst-case derivation programme is still OPEN: the superpolynomial lower
bound on $\Wb_{\Rm}$ and the bridge
$\mathrm{SAT}\in\mathrm P\Rightarrow \Wb_{\Rm}(N)\le\poly(N)$ (called B3 in
root \texttt{main.tex}). This run proves no new lower bound. It does three
things:
\begin{enumerate}
\item It records, with credit, the propositional bookkeeping already in the
corpus (LSPV-1/398, REFEREE357), and checks it mechanically
(\textbf{COMPUTATIONAL}, 6 admissible truth rows plus one hostile control).
\item It proves a \emph{strength calibration} (\textbf{PROVED/EXACT}, relative to
the corpus's paid ER2 $\subseteq_p \Rm$ inclusion): a superpolynomial
$\Wb_{\Rm}$ would in particular give superpolynomial lower bounds for
Extended Resolution and for \emph{unrestricted} Res($\oplus$), and the second
is reported open in the 2025 literature. So priority 1 of the directive is
at least as hard as a named open problem in proof complexity, and no
currently known hard family can supply it.
\item It assembles a hostile-family survival matrix (\textbf{KNOWN THEOREM}
per cited cell) and names the next concrete target.
\end{enumerate}

\section{Objects and exact quantifiers}
Inputs are canonical signed 3-CNFs with encoded length $|\langle F\rangle|$.
For a fixed Cook--Reckhow refutation system $R$ with nonnegative cost measure
$\cost_R$,
\[
 W_R(F)=\min_{\pi:F\vdash_R\bot}\cost_R(\pi),\qquad
 \Wb_R(N)=\max_{\substack{F\in\mathrm{UNSAT}\\|\langle F\rangle|\le N}}W_R(F).
\]
Write $\mathrm{BND}(R)$ for ``$\exists k\,\forall N:\Wb_R(N)\le N^k+k$'' and
$\mathrm{LB}(R)=\neg\mathrm{BND}(R)$. The bridge is
\[
 \mathrm{B3}(R):\qquad \mathrm{SAT}\in\mathrm P\ \Longrightarrow\ \mathrm{BND}(R).
\]
The directive's free-optimizer principle is the definition of $W_R$: the
minimum is taken over all legal derivations and the work of finding one is
not charged. This is sound for lower bounds. A machine time lower bound
follows from $W_R$ only for machines whose run \emph{induces} an
$R$-derivation of no larger cost. Supplying that for every polynomial-time
SAT algorithm is exactly B3, so the free-optimizer principle moves the
difficulty into B3 and does not remove it.

$S\le_p R$ means $R$ p-simulates $S$ on 3-CNF refutations with cost: there is
a polynomial $q$ such that every $S$-refutation of cost $c$ of $F$ gives an
$R$-refutation of $F$ with cost at most $q(|\langle F\rangle|+c)$.

\section{Bookkeeping identities (credited; not new)}
\begin{theorem}[PROVED/EXACT, elementary]\label{thm:book}
For all fixed $R$ and all $S\le_p R$:
\begin{enumerate}
\item[(a)] $\mathrm{B3}(R)\iff(\mathrm P\ne\mathrm{NP})\lor\mathrm{BND}(R)$.
\item[(b)] If $\mathrm{LB}(R)$ is a theorem, then $\mathrm{B3}(R)\iff\mathrm P\ne\mathrm{NP}$.
\item[(c)] $\mathrm{B3}(R)\land\mathrm{LB}(R)\Rightarrow\mathrm P\ne\mathrm{NP}$.
\item[(d)] $\mathrm{LB}(R)\Rightarrow\mathrm{LB}(S)$, and $\mathrm{B3}(S)\Rightarrow\mathrm{B3}(R)$.
\item[(e)] $\mathrm{B3}(R)\Rightarrow\mathrm{B3}(S)$ is \emph{not} a propositional consequence (hostile control).
\end{enumerate}
\end{theorem}
\begin{proof}
(a)--(c) are propositional. For (d), a cost-$c$ $S$-refutation yields an
$R$-refutation of cost $q(N+c)$, so $\mathrm{BND}(S)\Rightarrow\mathrm{BND}(R)$;
contrapose for the first half and compose with the antecedent for the
second. For (e), the assignment $\mathrm P=\mathrm{NP}$,
$\mathrm{BND}(R)$, $\neg\mathrm{BND}(S)$ is consistent with $S\le_p R$.
\end{proof}
Part (b) for Resolution is LSPV-1/398 (``the conditional is equivalent to
$\mathrm{SAT}\notin\mathrm P$''). The p-optimality reformulation under
$\mathrm{SAT}\in\mathrm P$ is REFEREE357. The script
\path{research/agents/claude/wcderiv/b3_logic_check.py} enumerates every
admissible truth row, all six admissible rows pass, and the control in (e)
fails as it should, on the row $(\mathrm P{=}\mathrm{NP},\mathrm{BND}(R),\neg\mathrm{BND}(S))$.
Raw output: \path{research/agents/claude/wcderiv/b3_logic_check_results.json}.
The run takes under one second, so no chunking was needed. No TeX engine is installed in this environment, so the source was checked for balanced environments by script and not compiled.

\emph{Consequence for system choice.} By (b), B3 is a weaker-than-target
bridge only for systems whose lower bound is \emph{not} already known. For
Resolution, Cutting Planes, Polynomial Calculus and bounded-depth Frege,
proving B3 is literally proving P$\ne$NP. So the directive's ``sufficiently
powerful'' fixed system has to sit above every system with a known
superpolynomial lower bound. $\Rm$ does, and that is why it is the correct
choice. The same fact makes priority 1 hard, as the next section shows.

\section{Strength calibration of priority 1}
\begin{lemma}[KNOWN THEOREM / corpus]\label{lem:incl}
(i) Extended Resolution (with Tseitin extension, paid per definition) satisfies
$\mathrm{ER}\le_p\Rm$ with node cost. This is the paid ER2-to-$\Rm$ inclusion
used as PROVED in LSPV-1/398 (root \texttt{main.tex}, \S\ LSPV-1/398).
(ii) [Cook 1975; Krajíček--Pudlák 1989] EF, and hence ER on CNF inputs,
p-simulates every Cook--Reckhow system $Q$ whose reflection principle has
polynomial-size EF proofs. In particular it p-simulates every system whose
soundness is provable in $\mathrm{PV}$/$S^1_2$. This covers Resolution, Cutting
Planes [Goerdt 1991 for Frege], Polynomial Calculus over finite fields with
Boolean axioms, bounded-depth Frege, Frege and Res($\oplus$)/Res(lin).
\end{lemma}

\begin{theorem}[PROVED/EXACT relative to Lemma~\ref{lem:incl}]\label{thm:cal}
If $\Wb_{\Rm}$ is superpolynomial along a cofinal sequence of 3-CNF lengths,
then on the same UNSAT 3-CNF family (i) minimum ER refutation size is
superpolynomial, and (ii) minimum refutation size in every system $Q$
covered by Lemma~\ref{lem:incl}(ii) is superpolynomial. In particular this
holds for unrestricted Res($\oplus$) and for Frege.
\end{theorem}
\begin{proof}
Apply Theorem~\ref{thm:book}(d) to $S=\mathrm{ER}\le_p\Rm$. Then compose
$Q\le_p\mathrm{ER}\le_p\Rm$, using transitivity of $\le_p$ (composition of
polynomials) on the identical signed source.
\end{proof}

\paragraph{Literature status of the consequences (preflight search, this run).}
A superpolynomial lower bound for general (unrestricted) Frege or EF is a
long-standing open problem. For unrestricted Res($\oplus$), a 2025 ECCC report
states that no superpolynomial size lower bound is known. Recent progress
is depth-restricted: Efremenko--Itsykson (depth $\Omega(n\log n)$) and
Bhattacharya--Chattopadhyay (arXiv:2507.23008, size
$\exp(\tilde\Omega(N^{\epsilon}))$ for depth $O(N^{2-\epsilon})$). Hence:

\begin{corollary}[PROVED/EXACT calibration]
Priority 1 of the shared directive, a superpolynomial $\Wb_{\Rm}$, implies
the currently open superpolynomial lower bound for unrestricted Res($\oplus$),
and also Frege and EF lower bounds. No hard family in the current
literature or corpus can supply it.
\end{corollary}

\begin{theorem}[CONDITIONAL equivalence]\label{thm:eq}
Assume (H1) the OPEN byte-domination lemma of REFEREE357, i.e.\ every
$\Rm$-refutation normalizes to one of bit length
$\poly(|\langle F\rangle|+\cost_{\Rm})$, and (H2) the soundness of the
normalized $\Rm$ verifier is PV-provable. Then $\Rm\equiv_p\mathrm{ER}$ on
3-CNF and
\[
 \mathrm{LB}(\Rm)\iff\mathrm{LB}(\mathrm{ER}),\qquad
 \mathrm B3(\Rm)\iff\bigl[\mathrm P=\mathrm{NP}\Rightarrow\mathrm{ER}\ \text{is p-bounded on 3-CNF}\bigr].
\]
\end{theorem}
\begin{proof}
(H1) and (H2) give $\Rm\le_p\mathrm{ER}$ by Lemma~\ref{lem:incl}(ii). Together
with Lemma~\ref{lem:incl}(i) this gives p-equivalence, and both identities
follow from Theorem~\ref{thm:book}(d) applied in both directions.
\end{proof}
Without (H1), the BLOB example of REFEREE357 shows that node cost can
undercount bytes, so the conditional cannot be dropped. Under (H1)--(H2), the
programme is exactly Cook's programme for EF: prove EF is not polynomially
bounded and prove that $\mathrm P=\mathrm{NP}$ forces EF to be polynomially
bounded. The only known route to the second is the Cook/Krajíček--Pudlák
reflection route, which needs a polynomial-time SAT algorithm whose
rejection-soundness is \emph{feasibly provable}. That is the
GODEL-REFLECTION357 firewall already in the corpus. $\mathrm{SAT}\in\mathrm P$
alone is not known to supply this.

\section{Hostile-family survival matrix (priority 4)}
U = known polynomial (or quasipolynomial, marked q) upper bound; L = known
superpolynomial lower bound; ? = open. The citations are the standard
attributions. Cells marked $\dagger$ were not re-read in primary form this run.

\begin{center}\small
\begin{tabular}{lccccccc}
\toprule
family & Res & PC/$\mathbb F_2$ & Res($\oplus$) & CP & AC$^0$-Frege & Frege & ER/$\Rm$\\
\midrule
PHP$^{n+1}_n$ & L [Haken] & L [Razborov] & U [Raz--Tzameret] & U [CCT87] & L [PBI, KPW] & U [Buss87] & U [Cook76]\\
odd Tseitin on expanders & L [Urquhart, BSW] & U (Gauss) & U & U$^{q}$ [DT20]$^\dagger$ & L [BS02/PBI]$^\dagger$ & U & U ($4h{-}1$, LSPV-2/341)\\
random 3-XOR & L [BSW] & U & U & ? & ? & U & U\\
random 3-CNF, $\Theta(n)$ cl. & L [CS88] & L [BSI, AR01] & ? (open) & L [HP17, FPPR17] & L [Gryaznov--Talebanfard 2025]$^\dagger$ & ? & ?\\
clique--colouring & L$^\dagger$ & ? & ? & L [Pudlák97] & L$^\dagger$ & ? & ?\\
\bottomrule
\end{tabular}
\end{center}

\paragraph{Reading.} Every row with a known lower bound in some column has a
known polynomial upper bound in a column at or below the ER/$\Rm$ column,
except random 3-CNF and clique--colouring, where both directions are open
from Res($\oplus$) or Frege upward. The fixed calculus $\Rm$ contains GF(2) row
additions and paid definitions, so PHP, Tseitin and random XOR are
\textbf{FALSIFIED/RETIRED} as hard candidates for $\Rm$. This is consistent
with the corpus: LSPV-1/398 (PHP), LSPV-2/341 (Tseitin $4h-1$). Random 3-CNF
near threshold is the only standard family that survives every system with
a known lower bound, is not known easy for Res($\oplus$), and keeps the
SAT/UNSAT boundary sharp. It is therefore the right candidate. By
Theorem~\ref{thm:cal}, though, establishing its hardness for $\Rm$ would
first require the open Res($\oplus$) and Frege lower bounds.

\section{Next target (single highest value)}
\textbf{T-RESXOR (OPEN):} prove a superpolynomial lower bound for
\emph{unrestricted} Res($\oplus$) on random 3-CNF, or on some explicit 3-CNF
family that is polynomially encoded. This is the weakest named milestone
implied by priority 1 (Theorem~\ref{thm:cal}). It lies inside the corpus's
existing RES+XOR fragment of $\Rm$, and current depth-restricted techniques
(lifting with stifling gadgets, depth-sensitive random restrictions) attack
it directly. A falsification-first companion: look for a polynomial-size
Res($\oplus$) or $\Rm$ refutation of random 3-CNF below the known CP and
bounded-depth Frege lower-bound regimes. Such a refutation would retire
random 3-CNF too.

\section*{Barrier-transfer preflight}
Main obstacles: (O1) any lower bound for $\Rm$ is at least an EF/Res($\oplus$)
lower bound; (O2) B3 needs feasible correctness of an arbitrary P-time
decider.
\begin{itemize}
\item \emph{Directly transferable theorem:} p-simulation transitivity and
Cook/Krajíček--Pudlák reflection simulation (used above).
\item \emph{Structural analogy needing a new theorem:} for O1, communication
lifting from parity decision trees (distributed computing cut lower bounds)
applies to Res($\oplus$) at bounded depth. The missing theorem is a
depth-free simulation of Res($\oplus$) by parity decision DAGs with a
size-preserving cut argument.
\item \emph{Failed transfer:} Algebraic Schur-square or rank arguments (PQ
items 1--3) bound GF(2)-linear reasoning, which $\Rm$ performs at linear cost
(Tseitin row). They are FAILED TRANSFER for $\Rm$-hardness, and they flag a
weakness: hidden linear structure is exploitable.
\end{itemize}
PQ ledger: item 5 (worst-to-average) is an ANALOGY/CANDIDATE only. Random
3-CNF is an average-case family, and no worst-case-to-average-case reduction
for coNP proof size is known. Item 6 (collision extraction) is
ANALOGY/CANDIDATE for the byte-domination lemma (H1). Items 4 and 7 are not
used.

\section*{CORPUS REVIEWED}
Repository tree (head \texttt{17cb1a45}); \texttt{AGENTS.md}, including the
2026-11-09 knot/link handoff (not engaged this run) and the world ledger
note; root \texttt{main.tex} B3/$\Wb_R$ material: lines on LSPV-1/350 MUS
quotient, /351 deficiency filter, /353 universal reflection, /357 Gödel
reflection firewall, /361 ER mask, /398 Resolution backtranslation firewall,
REFEREE357 p-optimality and byte domination, LSPV-2/341 Tseitin $4h{-}1$;
codex 2026-11-11 transfer recovery report; earlier Claude world ledger and
T-PORT notes. The newest activity at head consists of LSPV-3/581 and
LSPV-4/554 commits, which are all present and read back.

\section*{DUPLICATION / CONTRADICTION AUDIT}
Theorem~\ref{thm:book}(a)--(c) duplicates LSPV-1/398 and REFEREE357 and is
included only for completeness, credited. The direction
$\mathrm{LB}(\Rm)\Rightarrow\mathrm{LB}(\mathrm{ER})$ is implicit in the
corpus's inclusion lemma. The explicit chain to \emph{unrestricted
Res($\oplus$)} and the 2025 literature status check are new to the corpus as
far as repository search shows (no \texttt{main.tex} hit for Res($\oplus$)
lower-bound status). The survival matrix aggregates known results and claims
no novelty. No contradiction with any corpus claim was found.

\section*{WHAT CHANGED RELATIVE TO ROOT MAIN}
Nothing in root \texttt{main.tex} was edited. This agent owns no marked
region, and AGENTS.md reserves whole-document reconciliation for LSPV-5.
This report and the living log
\path{research/agents/claude/WORSTCASE_DERIVATION_LOG.md} are offered to
LSPV-5 as evidence-bearing input.

\section*{ADOPTED / REJECTED / RETIRED PRIOR COMPONENTS}
Adopted: $\Rm$ as the fixed system, ER2 inclusion, REFEREE357 (H1) as OPEN.
Retired for $\Rm$-hardness (reconfirmed, not newly retired): PHP, Tseitin,
random XOR. Rejected: reading the free-optimizer principle as removing the
bridge obligation.

\section*{STALE OR OVERSTATED CLAIMS FOUND}
None found in the read sections. One phrasing caution: ``sufficiently
powerful fixed system'' should be read as ``above every system with a known
lower bound'' (Theorem~\ref{thm:book}(b)). Otherwise B3 is not weaker than
P$\ne$NP.

\section*{Predecessor / write check}
This agent is outside the LSPV-1..5 cyclic stagger. The latest scheduled
peer activity at head (LSPV-3/581, LSPV-4/554) has committed and verified
artifacts, so no recovery prompt was needed.

\section*{References}
\begin{itemize}\small
\item S.~Cook, Feasibly constructive proofs and the propositional calculus, STOC 1975.
\item S.~Cook, R.~Reckhow, The relative efficiency of propositional proof systems, JSL 44, 1979.
\item J.~Krajíček, P.~Pudlák, Propositional proof systems, the consistency of first order theories and the complexity of computations, JSL 54, 1989.
\item A.~Haken, The intractability of resolution, TCS 39, 1985. A.~Urquhart, Hard examples for resolution, JACM 1987.
\item E.~Ben-Sasson, A.~Wigderson, Short proofs are narrow, JACM 2001. V.~Chvátal, E.~Szemerédi, Many hard examples for resolution, JACM 1988.
\item E.~Ben-Sasson, R.~Impagliazzo, Random CNFs are hard for the polynomial calculus, FOCS 1999; M.~Alekhnovich, A.~Razborov, Lower bounds for polynomial calculus: non-binomial case, FOCS 2001.
\item P.~Hrubeš, P.~Pudlák, Random formulas, monotone circuits, and interpolation, FOCS 2017; N.~Fleming, D.~Pankratov, T.~Pitassi, R.~Robere, Random $\Theta(\log n)$-CNFs are hard for cutting planes, FOCS 2017.
\item R.~Raz, I.~Tzameret, Resolution over linear equations and multilinear proofs, APAL 2008.
\item Y.~Efremenko, D.~Itsykson, lower bounds for regular and bounded-depth Res($\oplus$), 2024--2025; S.~Bhattacharya, A.~Chattopadhyay, arXiv:2507.23008 (2025); ECCC 2025 reports cited in the preflight search.
\item S.~Buss, Polynomial size proofs of the propositional pigeonhole principle, JSL 1987; W.~Cook, C.~Coullard, Gy.~Turán, 1987; P.~Pudlák, Lower bounds for resolution and cutting plane proofs and monotone computations, JSL 1997.
\end{itemize}
\textbf{P versus NP remains OPEN.}
\end{document}
% END_OF_VERIFIED_REPORT::CLAUDE::WCDERIV-0001
